\documentclass[11pt]{article}
\usepackage[a4paper,margin=27mm]{geometry}
\usepackage{amsmath,amssymb,amsthm,mathtools}
\usepackage[T1]{fontenc}
\usepackage{lmodern,microtype,xurl}
\usepackage[colorlinks=true,linkcolor=blue,citecolor=blue,urlcolor=blue]{hyperref}
\newtheorem{theorem}{Theorem}
\newtheorem{lemma}{Lemma}
\newcommand{\F}{\mathbb F}
\title{Ordinary Polynomial Multiplication over Fixed Finite Fields\\
  in $O(n\log n)$ Time\\
  \large A baseline companion and fixed-basis reduction}
\author{\'Arp\'ad Csizmadia\\
  \small\href{mailto:arpad.csizmadia.public@gmail.com}{\texttt{arpad.csizmadia.public@gmail.com}}}
\date{31 October 2026\\ \small Version 0.1.0-preprint. Apache-2.0.}
\hypersetup{pdfauthor={\'Arp\'ad Csizmadia},pdftitle={Ordinary Polynomial Multiplication over Fixed Finite Fields in O(n log n) Time}}
\begin{document}
\maketitle

\begin{abstract}
For every individually fixed finite field $\F_q$, we derive ordinary
dense polynomial multiplication in $O_q(n\log(n+2))$ deterministic bit
time on a fixed finite-alphabet Turing machine with a fixed finite
number of sequential tapes. The characteristic-two branch uses the
accompanying Pascal-transform binary multiplier; the odd-prime branch
weakens the accompanying exponent $9999/10000$. A fixed-basis expansion
reduces a degree-$e$ extension to $e^2$ prime-field products and
$O_q(n)$ encoding and recombination work, including ordinary coefficient
input/output. This note records a baseline consequence of shared proof
sources. Complete Lean formalization is unfinished.
\end{abstract}

\paragraph{AI assistance and review status.}
This work was developed with substantial generative AI assistance in
mathematical exploration, proof development, writing, coding and source
review. OpenAI Codex was used for release preparation and source
inspection; the platform reference~\cite{CodexPlatform} identifies the
tool, not a mathematical premise. A complete historical model/version
log is not available. Full Lean formalization is unfinished and planned
for a later update, with no completion date promised. This preprint is
open for mathematical review; independent human proof review is not claimed.

\section{Scope, prior work and the role of this note}
The question is the bit complexity of multiplying ordinary dense
univariate polynomials over an individually fixed finite field.
Harvey, van der Hoeven and Lecerf established a F\"urer-type bound
with the factor $8^{\log^* n}$ for each fixed prime
field~\cite{HVHL17}. Harvey and van der Hoeven then obtained the
unconditional fixed-prime bound
\[
 O_p\!\left(n\log n\,4^{\log^* n}\right)
\]
in the deterministic multitape Turing model~\cite[Theorem 1.1]{HVH19Cyclo}.
For a fixed extension degree the same asymptotic form follows by
coordinate reduction. These are important predecessors for the
fixed-field comparison in this note.

Their later finite-field theorem gives
$O(n\log q\log(n\log q))$ uniformly in $q$, under the hypothesis
that a Linnik constant $L<1+2^{-1162}$ exists
\cite[Theorem 1.2]{HVH22Finite}. Thus its fixed-$q$ specialization
is a conditional $O_q(n\log n)$ theorem. Its quantifiers and
number-theoretic hypothesis should remain explicit. Harvey's
October 2026 discussion separately identifies fixed-field polynomial
multiplication as a problem not directly settled by the new integer
multiplication claim, and records the unconditional cyclotomic
benchmark~\cite[Section 4, p.~4]{Harvey26}. That discussion supplies
context, not verification or endorsement of this note.

This history explains the role of the $n\log n$ benchmark: the
unconditional fixed-field algorithms reached near-linear time,
while the sharper uniform theorem retained its explicit hypothesis.
The accompanying new constructions claim to close the remaining
fixed-field gap. This note makes their shared baseline consequence
and the elementary extension-field reduction easy to identify.

The present baseline combines two arguments in the accompanying
preprints. For characteristic two it uses the Pascal-transform route
of~\cite{SimpleBinary}. For odd characteristic it uses only the
odd-field proposition of~\cite{GeneralBest}. It does not use that
paper's stronger combined all-field theorem or the separate
controlled-CRT binary multiplier. The fixed-basis reduction below
makes the resulting field quantifiers and ordinary input/output
boundary explicit. This is a companion consequence of shared proof
sources, not a separate independent discovery or a practical speed claim.

\section{The fixed-field statement}
Let $q=p^e$ be fixed, with $p$ prime and $e\ge1$. Choose a fixed
$\F_p$-basis $b_0,\ldots,b_{e-1}$ of $\F_q$. Represent a field
element by its $e$ prime-field coordinates, each in a fixed binary
encoding. The field, basis and structure constants are fixed before
the input length varies; they are not growing inputs or advice
depending on $n$.

An input is the two ascending coefficient lists of polynomials
$A,B\in\F_q[X]$ of degree less than $n$, with $n\ge1$, including
leading zero coefficients. The required output is the ascending list
of all $2n-1$ coefficients of $AB$, also retaining leading zero
coefficients. Input and output have $\Theta_q(n)$ bits. Time includes
all encoding, setup, movement, component calls, decoding and output.

\begin{theorem}[Ordinary fixed-field baseline]
\label{thm:baseline}
For every individually fixed finite field $\F_q$, the proof routes
in~\cite{SimpleBinary,GeneralBest} give a deterministic algorithm
for this ordinary product in
\[
 O_q\!\left(n\log(n+2)\right)
\]
bit steps. The machine has a fixed finite alphabet and a fixed finite
number of sequential one-dimensional tapes. Its finite control,
tape count and implied constants may depend on $q$, but not on $n$.
\end{theorem}

\section{Prime-field inputs and inherited dependencies}
Write $M_p(n)$ for the cost of an ordinary prime-field multiplier
with the input/output convention just stated.

For $p=2$, the completed reduction and binary tensor construction
in~\cite{SimpleBinary} state
$M_2(n)=O(n\log(n+2))$. That route includes its own evaluation,
pointwise product and residue decoder. It retains the arbitrary-width
equal-chunk tape-interchange result from OpenAI's
preprint~\cite[Lemma 12, Section 4]{OpenAI26}. This note inherits
that cited theorem dependency; it does not infer a binary polynomial
multiplier simply from an integer-multiplication bound.

For an odd prime $p$, Proposition~1 of~\cite{GeneralBest}, specialized
to extension degree one, states
\[
 M_p(n)=O_p\!\left(n(\log(n+2))^{9999/10000}\right).
\]
Its bundled odd-field proof uses the unconditional cyclotomic host
of~\cite{HVH19Cyclo}, the accompanying native constructions and a
named stage-two interchange theorem~\cite{Stage2}. The proof sources
and exact movement adapter are identified in that companion's
\path{proofs/PROOF_SOURCES.md}. This odd proposition is proved
separately from its characteristic-two combination.

Since $\log_2(n+2)>1$ for $n\ge1$ and $9999/10000<1$, weakening
this bound gives $M_p(n)=O_p(n\log(n+2))$. In particular, the odd
part of this baseline does not introduce a new Linnik-type
hypothesis. The present note relies on the written cited arguments;
it does not independently certify their correctness or formalize
their full machine constructions.

\section{Fixed-basis reduction and tape costs}
\begin{lemma}[Constant-degree coordinate reduction]
\label{lem:extension}
For a fixed extension $\F_q/\F_p$ of degree $e$, an ordinary
prime-field multiplier of cost $M_p(n)$ yields an ordinary
$\F_q$ multiplier of cost
\[
 O_q\!\left(e^2 M_p(n)+n\right).
\]
Here $e$ is fixed; displaying it records the number of component
products, not a uniform growing-extension complexity claim.
\end{lemma}
\begin{proof}
Expand the coefficient lists in the fixed basis to write
\[
 A=\sum_{i=0}^{e-1} b_i A_i,\qquad
 B=\sum_{j=0}^{e-1} b_j B_j,
 \qquad A_i,B_j\in\F_p[X],\quad \deg A_i,\deg B_j<n.
\]
Each $A_i$ and $B_j$ is an ordinary coefficient list in ascending
degree order. Fixed-width coefficient decoding and a fixed number
of sequential scans write these $2e$ lists in $O_q(n)$ steps.
An initial length scan and all component-call descriptors have the
same bound. Any other fixed field-element encoding can be converted
by fixed finite tables at this cost.

Compute the $e^2$ ordinary products
\[
 C_{ij}=A_iB_j\quad(0\le i,j<e)
\]
with the prime-field multiplier. Each result has $2n-1$ ordinary
coefficients. Calls and input copies are sequential, and the fixed
number of component machines can be incorporated into one finite
control with a fixed finite set of work tapes. Their total cost,
including their stated setup and cleanup, is
$e^2M_p(n)+O_q(n)$.

The fixed structure constants $\mu_{ijr}\in\F_p$ satisfy
\[
 b_i b_j=\sum_{r=0}^{e-1}\mu_{ijr}b_r.
\]
Consequently, if $D_r=\sum_{i,j}\mu_{ijr}C_{ij}$, then
$AB=\sum_r b_rD_r$. Recombination uses a fixed number of scans
over the component lists: the $k$-th output is formed only from
their $k$-th coefficients. All operations on $\F_p$ and all fixed
scalar multiplications have constant bit cost depending on $q$.
At most $O(e^3n)=O_q(n)$ such work is required. A fixed number
of tapes retains the component streams, so these scans assume no
unit-cost random access or free global data permutation.

Finally encode $\sum_r b_r(D_r)_k$ and write it for
$k=0,\ldots,2n-2$. This is precisely the ordinary product list;
no cyclic wrap, transformed output or missing high coefficients
remain. Output writing and disposal of temporary streams cost
$O_q(n)$, proving the stated bound.
\end{proof}

\begin{proof}[Proof of Theorem~\ref{thm:baseline}]
Section~3 supplies $M_p(n)=O_p(n\log(n+2))$ for each fixed prime
$p$, with the cited dependency chain appropriate to its characteristic.
Apply Lemma~\ref{lem:extension}. Since $p,e$ and all field data
are fixed, $e^2$ is constant and the linear codec/recombination term
fits the stated bound. The finitely many small lengths can be
handled by ordinary finite-table arithmetic on the same machine.
\end{proof}

\section{Verification and relation to stronger claims}
This note isolates the general $n\log n$ baseline and its fixed-basis
proof. The stronger general companion~\cite{GeneralBest} states
a strict logarithmic saving by a different combination of its
odd-field proposition with the principal binary paper. The baseline
proved here instead uses the accessible binary companion. The two
general statements share the odd-field argument and should be
credited and reviewed with that shared dependency visible.

No new finite experiment or end-to-end Lean verification is offered
by this reduction note. Compilation and package hashes, when supplied,
associate the document with its source; they do not certify the
asymptotic theorem. Full Lean formalization remains unfinished.
The finite-field quantifiers, ordinary coefficient representation
and tape costs are part of any future statement-level verification.

\begingroup
\small
\raggedright
\begin{thebibliography}{9}
\bibitem{HVHL17}
David Harvey, Joris van der Hoeven and Gr\'egoire Lecerf.
\newblock Faster polynomial multiplication over finite fields.
\newblock \emph{Journal of the ACM}, 63(6), Article 52, 2017.
\newblock \href{https://doi.org/10.1145/3005344}{doi:10.1145/3005344}.

\bibitem{HVH19Cyclo}
David Harvey and Joris van der Hoeven.
\newblock Faster polynomial multiplication over finite fields using
cyclotomic coefficient rings.
\newblock \emph{Journal of Complexity}, 54, Article 101404, 2019.
\newblock \href{https://doi.org/10.1016/j.jco.2019.03.004}{doi:10.1016/j.jco.2019.03.004}.
\newblock \href{https://www.texmacs.org/joris/cyclomult/cyclomult.pdf}{Author manuscript}.

\bibitem{HVH22Finite}
David Harvey and Joris van der Hoeven.
\newblock Polynomial multiplication over finite fields in time $O(n\log n)$.
\newblock \emph{Journal of the ACM}, 69(2), Article 12, 1--40, 2022.
\newblock \href{https://doi.org/10.1145/3505584}{doi:10.1145/3505584}.
\newblock Theorem 1.2 is conditional on the stated Linnik hypothesis.
\href{https://www.texmacs.org/joris/ffnlogn/ffnlogn.pdf}{Author manuscript}.

\bibitem{Harvey26}
David Harvey.
\newblock On the recently proposed $o(n\log n)$ integer multiplication algorithm.
\newblock Unpublished note, 20 October 2026, Section 4, p.~4.
\newblock \url{https://web.maths.unsw.edu.au/~davidharvey/research/subnlogn.pdf}.

\bibitem{SimpleBinary}
\'Arp\'ad Csizmadia.
\newblock \emph{Exact Binary Polynomial Multiplication from Fast Pascal Transforms}.
\newblock Accompanying AI-assisted research preprint, version 0.1.0-preprint,
2026; \path{papers/02-simple-binary/} in this collection.

\bibitem{GeneralBest}
\'Arp\'ad Csizmadia.
\newblock \emph{Ordinary Polynomial Multiplication over Every Fixed Finite Field}.
\newblock Accompanying AI-assisted research preprint, version 0.1.0-preprint,
2026; Proposition 1 and its bundled proof sources,
\path{papers/01-general/} in this collection.

\bibitem{OpenAI26}
OpenAI.
\newblock \emph{Integer multiplication below $n\log n$}.
\newblock Research preprint, 23 September 2026; Lemma 12, Section 4.
\newblock Citation revision \texttt{4eaee87548250b2682affda7d05d947c05554bae}.
\newblock \url{https://github.com/openai/math/blob/4eaee87548250b2682affda7d05d947c05554bae/preprints/Integer-multiplication-below-n-log-n-September-23-2026/build/sections/04-swap.tex}.

\bibitem{Stage2}
\texttt{platypii/integer-mult-bounds-lean} contributors.
\newblock Stage-two equal-width interchange bound, declaration
\path{IntegerMultBounds.MathlibTM2.M8.Stage2Machine.Bound.uniform_bound},
commit \texttt{fefd699f9389af961046997796261333539e5909}.
\newblock \href{https://github.com/platypii/integer-mult-bounds-lean/blob/fefd699f9389af961046997796261333539e5909/IntegerMultBounds/MathlibTM2/M8/Stage2Machine/Bound.lean}{Pinned theorem source}.

\bibitem{CodexPlatform}
Nicolas Bonamy and Derrick Choi.
\newblock Codex as a platform: build on the open agent harness.
\newblock OpenAI Developers, 19 August 2026.
\newblock \url{https://developers.openai.com/blog/codex-as-a-platform}.
\end{thebibliography}
\endgroup
\end{document}
